\medskip\noindent\textbf{One cone conjugator and all five exact filling equations.}
The kernel equation (F.7.44) has an exact formulation
using a single element of $H$.  The full profinite
generation assertion needed in this formulation
follows from the actual degree-one cusp of a finite
orbifold cover.

First, every closed subgroup $L\leq E$ satisfying
\[
              d\in L,\qquad d^6\in\overline{[L,L]}
                                                               \tag{F.7.49}
\]
is $E$.  Indeed a proper such $L$ is contained in a
proper open subgroup $U$: separate a point outside
$L$ in a finite quotient and take the inverse image
of the image of $L$.  The subgroup
$U\cap\Delta$ has the same finite index, and its full
completion is $U$.  In the corresponding connected
orbifold cover, $d\in U\cap\Delta$ gives a width-one
cusp through the base point.  A proper cover has
at least one more cusp, since the remaining cosets
cannot lie in that length-one cycle of $d$.
Its cone orders divide two or three.
For any prime $p>3$, the cone classes vanish in
$H_1(U,\F_p)$; the cusp classes have precisely their
total-sum relation, together with independent
handle classes.  Each cusp class is therefore
nonzero.  The class of $d^6$ is six times the
specified width-one cusp and is nonzero, contradicting
(F.7.49).

It follows that every continuous endomorphism
$\rho:E\to E$ satisfying
$\rho(d)=t d^\mu t^{-1}$ for a full unit
$\mu\in\widehat\Z^\times$ is an automorphism.
The compact image $L=t^{-1}\rho(E)t$ contains
$d^\mu$ and hence its $\mu^{-1}$ power $d$.
Since $E^{\rm ab}=C_2\times C_3$, the element
$d^6$ belongs to the closed commutator subgroup
of $E$.  Thus $d^{6\mu}\in\overline{[L,L]}$,
and its $\mu^{-1}$ power gives (F.7.49).
Surjectivity follows, and finite generation of $E$
gives Hopficity.  Both the conjugation and the
unit-root operation here are performed on the
specified image subgroup; no pro-$p$ quotient
replaces $E$.

Define the following compact solution space:
\[
 \mathscr S=
 \{\,g\in H:(d g b^{-1}g^{-1})^2=1\,\},\qquad
 \eta_g(d)=d,\quad
 \eta_g(b)=g b g^{-1},\quad
 \eta_g(a)=d g b^{-1}g^{-1}.
                                                               \tag{F.7.50}
\]
The equation gives the $a^2$ relation, and the
$b^3$ relation holds by conjugation.  Hence the
displayed substitutions define an endomorphism
of $E$, and the preceding full boundary test makes
it an automorphism.  Its action on $A_5$ is the
identity because $g\in H$.  Consequently its outer
class is the element of $GT_0^1$ specified by
(F.7.23).

Every element of that $GT_0^1$ has this form.
For $\gamma=(1,f)$, use the actual
$\eta_j=\operatorname{Inn}(d^{-j})\eta_f$ in
(F.7.27).  The full free profinite product torsion
theorem makes $\eta_j(b)$ conjugate to $b$ or
$b^{-1}$; these are the only order-three factor
elements in $C_2\amalg C_3$
\cite[Theorem A and Section 3]{A-HerfortRibes1985}.
Fixing $d$ forces the identity on
$E^{\rm ab}=C_6$, which excludes $b^{-1}$.
Choose $u$ with $\eta_j(b)=ubu^{-1}$.
Since the quotient action is the identity,
$q_5(u)\in C_{A_5}(q_5(b))=\langle q_5(b)\rangle$.
Write $q_5(u)=q_5(b)^r$ and replace $u$ by
$g_f=ub^{-r}\in H$.  The same source's full
free-factor centralizer theorem gives
$C_E(b)=\langle b\rangle$
\cite[Theorem B and Section 4]{A-HerfortRibes1985}.
As $H\cap\langle b\rangle=1$, this $g_f$ is unique.
The conjugation map $H\to E$, $g\mapsto gbg^{-1}$,
is an embedding of a compact space and thus has
a continuous inverse on its image.
We have the exact identities
\[
 \eta_j=\eta_{g_f},\qquad
 h_1=g_f b g_f^{-1}b^{-1},\qquad
 f=b d^j g_f b^{-1}g_f^{-1}d^{-j}.
                                                               \tag{F.7.51}
\]
The last formula retains the particular finite
$j$ selected for the original normalized pair;
it follows from $\eta_f=\operatorname{Inn}(d^j)\eta_j$.

Two solutions $g,h\in\mathscr S$ have the same
outer class exactly when
\[
                   h=d^{5t}g\quad(t\in\widehat\Z).
 \qquad
 \mathscr S/\bigl(g\sim d^{5t}g\bigr)
                         \xrightarrow{\ \sim\ }GT_0^1.
                                                               \tag{F.7.52}
\]
To prove the assertion, an inner comparison of
$\eta_g,\eta_h$ must centralize $d$, so its
conjugator lies in $\overline{\langle d\rangle}$
by (F.7.24).  Both quotient actions are the
identity; since $Z(A_5)=1$, that conjugator is
$d^{5t}\in H$.  Comparing the two images of $b$
then gives $h=d^{5t}g$, because the possible
remaining factor lies in
$H\cap C_E(b)=1$.  Conversely left multiplication
by $d^{5t}$ changes $\eta_g$ by precisely
$\operatorname{Inn}(d^{5t})$ and preserves
(F.7.50).
This proves the bijection, and compactness gives
the stated homeomorphism of the solution-orbit
space with the compact group on the right.

For the first detector the exponent condition
sharpens to one procyclic condition:
\[
 \begin{split}
 \ker\mathcal D_{J_0}
   &=\{\, (1,f):W_f=x^{A_f}y^{-A_f}\,\}\\
   &=\{\, (1,f):p_0(g_f)\in
                            \overline{\langle x\rangle}\,\},\\
 \mathscr S\cap\ker(p_0|_H)
   &\xrightarrow{\ \sim\ }\ker\mathcal D_{J_0}.
 \end{split}                                                   \tag{F.7.53}
\]
For the proof let $\varphi_g$ be the actual
automorphism of $E^+$ descended from $\eta_g$
through $p_0$.  It sends $b$ to
$p_0(g)b p_0(g)^{-1}$, and restricts to the
surface automorphism in (F.7.43).
When its outer restriction is the identity,
that restriction is $\operatorname{Inn}(x^{A_f})$,
by (F.7.44).  Since $C_{E^+}(K)=1$, two
automorphisms of $E^+$ with this same restriction
are equal: their quotient fixes $K$ pointwise,
and its difference on each element centralizes
$K$.  It follows that
$\varphi_g=\operatorname{Inn}(x^{A_f})$.
Comparing the images of $b$ gives
$p_0(g)=x^{A_f}$, since
$K\cap C_{E^+}(b)=1$.
Therefore
$W_f=x^{A_f}b x^{-A_f}b^{-1}=x^{A_f}y^{-A_f}$.
Conversely this expression is exactly the inner
kernel equation (F.7.44).
Finally (F.7.52) replaces a solution with
$p_0(g)=x^t$ uniquely by $d^{-5t}g$ with
$p_0(d^{-5t}g)=1$.  This gives the last
homeomorphism in (F.7.53).
The group law on this slice is explicitly
$g\star h=\eta_g(h)g$: here $\eta_g$ descends
to the identity through $p_0$, so the product
still lies in the slice.

The other four representative fillings can be
written in the same actual coordinates.
An independent index-sixty Schreier calculation
gives $H$ with $61$ initial generators and
$120$ rewritten relators.  Fifty single-occurrence
Tietze eliminations leave eleven generators and
no relators.  For each of the five triples in
(F.7.32), the nine omitted cusp relators give
nine more such eliminations, leaving a full
free basis $P,Q$ with the following values:
\[
 \begin{array}{c|c|ccc|cc}
 J&m&c_0&c_m&c_{\rm third}&P\text{ in }x,y&Q\text{ in }x,y\\ \hline
 \{0,1,2\}&1&QP&Q^{-1}&P^{-1}&yx&y^{-1}\\
 \{0,1,4\}&1&Q&P^{-1}&Q^{-1}P&y^{-1}&x\\
 \{0,1,8\}&1&QP&Q^{-1}&P^{-1}&yx&y^{-1}\\
 \{0,1,10\}&1&Q&P^{-1}&Q^{-1}P&y^{-1}&x\\
 \{0,4,9\}&4&P&P^{-1}Q&Q^{-1}&x&xy
 \end{array}                                                   \tag{F.7.54}
\]
The last two columns are inverse substitutions:
they identify $c_0,c_m$ with the specified full
basis $x,y$ of $K$.
All omitted cusp words go to one.
Thus the resulting maps
$\pi_J:H\to K$ are precisely the actual fillings
(F.7.29) followed by these based quotient
isomorphisms.  The complete original generators,
their words in $a,b$, all relators, all
eliminations, and both inverse basis maps
are recorded in the certificate.
For $J_0$, this is exactly $p_0|_H$,
as both identifications send $c_0,c_1$ to $x,y$.

For $g\in\mathscr S$ put $h(g)=g b g^{-1}b^{-1}\in H$.
Use the original $y=e^2\in E$ in the following
source word:
\[
 \begin{array}{c|c}
 J& V_J(g)\\ \hline
 \{0,1,2\},\{0,1,4\},\{0,1,8\},\{0,1,10\}
       &\pi_J(h(g))\\
 \{0,4,9\}&
          \pi_J\bigl(h(g)y h(g)^{-1}y^{-1}\bigr).
 \end{array}                                                   \tag{F.7.55}
\]
The whole word in the last row lies in $H$,
because $H$ is normal; the map $\pi_J$ is
applied to that word and is not evaluated
separately on $y$ outside its source.
To check the row, the exact label-four
whisker is $(ba)^2=y$, so $c_4=yc_0y^{-1}$.
The substitutions (F.7.50) fix $d$ and $x$,
and give $\eta_g(y)=h(g)y h(g)^{-1}$.
Consequently the conjugator for $c_4$ is
$h(g)y h(g)^{-1}y^{-1}$.
For the first four rows the conjugator for
$c_1=e^5$ is $h(g)$, as in (F.7.51).
Every descended map therefore fixes $x$ and
sends $y$ to $V_J(g)yV_J(g)^{-1}$.
If $A_J=\ell_x(V_J(g))$ and
$B_J=\ell_y(V_J(g))$, all five exact kernel
equations are
\[
 \begin{split}
 \mathcal D_J([\eta_g])&=
       [\,x\mapsto x,\ 
          y\mapsto F_J(g)^{-1}yF_J(g)\,],\\
 F_J(g)&=y^{B_J}V_J(g)^{-1}x^{A_J},\\
 \mathcal D_J([\eta_g])=1
             &\quad\Longleftrightarrow\quad
                   V_J(g)=x^{A_J}y^{B_J}.
 \end{split}                                                   \tag{F.7.56}
\]
Here $[\eta_g]$ is read as its exact $GT_0^1$
class under (F.7.23).
This supplies the same single full profinite
$g$ in every row, together with the
quadratic equation (F.7.50).

The four $C_3$-stable triples of (F.7.34) admit
one further common word.  With the same
$\tau$ from (F.7.33), define
\[
 R(g)=g^{-1}\eta_g(\tau)g\tau^{-1}\in H.
 \qquad
 \mathcal D_{J_i}([\eta_g])=1
       \quad\Longleftrightarrow\quad p_i(R(g))=1
       \quad(0\leq i\leq3).
                                                               \tag{F.7.57}
\]
Indeed $r=b$ is conjugated by $g$, while
$s=\tau b\tau^{-1}$ is conjugated by
$v(g)=\eta_g(\tau)g\tau^{-1}$.
After filling, $r,s$ generate $E^+$.
An inner surface action extends uniquely across
$E^+$ and its conjugator must be $p_i(g)$,
by $C_{E^+}(K)=1$ and
$K\cap C_{E^+}(b)=1$.
Its agreement on $s$ is then exactly
$p_i(v(g))=p_i(g)$, because the other finite
factor also has trivial centralizer intersection
with $K$.  This is (F.7.57), in both directions.
Replacing $g$ by $d^{5t}g$ and $\eta_g$ by its
specified inner change leaves $R(g)$ unchanged,
so this common-word comparison retains the
outer normalization.

The actual geometric residual group in (F.7.25)
therefore maps into the simultaneous solution
set (F.7.50), (F.7.53), and (F.7.56), with the
common-word constraints (F.7.57).
Triviality of that solution set has not been
proved.  The nontrivial common source kernel
in (F.7.41) supplies no solution to the additional
quadratic equation.  Conversely the exact
section and the primewise images in (F.7.48)
supply no step replacing these full quotient
equations.

The executed standard-library program
\path{code/part_a/modular_five_fillings_f7.py}
records all $95$ Tietze eliminations, all five
inverse full-basis maps, every original cusp
word, the order-three finite centralizer,
the five $j$-recovery word identities in
(F.7.51), and the exact label-four conjugator.
The orbifold generation test, the full torsion
and centralizer inputs, and the simultaneous
solution-space identifications are proved above.
